% CV académico — plantilla mínima
% Compila con pdflatex sin paquetes exóticos. Funciona en Overleaf sin configurar nada.
% Reemplaza TODO lo que esté entre corchetes. Borra las secciones que no apliquen:
% una sección vacía resta más de lo que suma.

\documentclass[11pt,letterpaper]{article}

\usepackage[utf8]{inputenc}
\usepackage[T1]{fontenc}
\usepackage[margin=1in]{geometry}
\usepackage{titlesec}
\usepackage{enumitem}
\usepackage[hidelinks]{hyperref}
\usepackage{parskip}

% Secciones en versalitas con una regla debajo: convención académica sobria.
\titleformat{\section}{\large\scshape}{}{0em}{}[\vspace{-0.6em}\rule{\textwidth}{0.4pt}]
\titlespacing{\section}{0pt}{1.2em}{0.5em}

\setlist[itemize]{leftmargin=1.2em, itemsep=0.15em, topsep=0.25em}
\pagestyle{empty}

% Encabezado de entrada: título a la izquierda, fecha a la derecha.
\newcommand{\entry}[2]{\noindent\textbf{#1}\hfill #2\par}

\begin{document}

%========================= HEADER =========================
\begin{center}
  {\LARGE \textbf{[Nombre Completo]}}\\[0.4em]
  [correo@institucional.edu.pe] $\cdot$ [Ciudad, Perú]\\
  \href{[url-scholar]}{Google Scholar} $\cdot$
  \href{[url-github]}{GitHub} $\cdot$
  \href{[url-sitio]}{Sitio personal}
\end{center}

%========================= EDUCATION =========================
\section{Education}

\entry{[B.Sc. in Computer Science], [Universidad]}{[2021--2026]}
\begin{itemize}
  \item Cumulative average: [16.2]/20 ([Quinto Superior] --- top [20]\% of a cohort of [142])
  \item Thesis: \emph{[Título de la tesis]} --- Advisor: Prof.\ [Nombre]
  \item Relevant coursework: [Machine Learning, Distributed Systems, Numerical Methods]
\end{itemize}

%========================= PUBLICATIONS =========================
% Borra esta sección completa si no tienes publicaciones. Es normal no tenerlas
% al postular a una maestría; una sección vacía llama la atención hacia el vacío.
\section{Publications and Preprints}
\begin{enumerate}[leftmargin=1.4em]
  \item [Apellido, N.], \textbf{[Tu Apellido, N.]}, [Apellido, N.]
        ``[Título del trabajo].'' \emph{[Venue]}, [año].
        \href{[doi-o-repo]}{[DOI / repositorio]}
\end{enumerate}

%========================= RESEARCH =========================
\section{Research Experience}

\entry{[Undergraduate Researcher], [Laboratorio], [Universidad]}{[2025--2026]}
\begin{itemize}
  \item Investigated [problema]. Implemented [método] in [herramienta] and
        evaluated against [baseline] on [dataset, tamaño].
  \item Achieved [métrica], a [X]\% improvement over [baseline].
  \item Code and experiment logs: \href{[url]}{[repositorio]}
\end{itemize}

%========================= TECHNICAL =========================
% Tu experiencia de industria, reencuadrada: método, baseline, número.
% Nada de métricas de negocio sin traducir.
\section{Technical Experience}

\entry{[Software Engineer], [Empresa]}{[2024--2026]}
\begin{itemize}
  \item Built and profiled [sistema] serving [escala].
  \item Reduced [métrica técnica] from [X] to [Y] through [técnica específica].
  \item Designed the benchmarking harness used to validate the change.
\end{itemize}

%========================= SOFTWARE =========================
\section{Scientific Software and Open Source}
\begin{itemize}
  \item \href{[url]}{\textbf{[repo-name]}} --- [una línea sobre qué hace].
        [Lenguaje]. Includes a Dockerfile and a fixed-seed evaluation script.
\end{itemize}

%========================= AWARDS =========================
\section{Honors and Awards}
\begin{itemize}
  \item [Premio o beca], [otorgante], [año]
\end{itemize}

%========================= SKILLS =========================
\section{Technical Proficiencies}
\begin{itemize}
  \item \textbf{Languages:} [Python, C++, Rust, TypeScript]
  \item \textbf{Research stack:} [PyTorch, CUDA, NumPy, SciPy]
  \item \textbf{Tooling:} [Docker, Git, SLURM, \LaTeX]
\end{itemize}

%========================= LANGUAGES =========================
\section{Languages}
\begin{itemize}
  \item Spanish (native) $\cdot$ English ([C1] --- [IELTS 7.5, 2026])
\end{itemize}

\end{document}
